\begin{abstract}\RaggedRight
Security alert grouping is risky when multiple apparent corroborating claims descend from one observation, arrive after the decision time, or disagree about event identity. We examine a constrained correlation, verification, and prioritization pipeline that forms initial alert fragments, accepts cross-fragment relationships only under explicit provenance and consistency checks, and ranks the resulting groups. A narrowly scoped endpoint fallback substitutes original endpoint process context for a missing process claim only when timely, independent authentication and network observations support the same relationship. We compare this frozen candidate against a matched frozen baseline on an untouched, controlled synthetic holdout of 15 queues and 5,018 eligible alerts. Incident Recall@20 rose from 250/300 (83.33\%) to 260/300 (86.67\%); group recoverability rose from 263/300 (87.67\%) to 281/300 (93.67\%); and P@10 rose from 107/150 (71.33\%) to 113/150 (75.33\%). Both pipelines had zero false grouped pairs and zero mixed groups under simulator truth. The candidate had missed its development gate by one incident, and the subsequent decision to evaluate it unchanged on fresh holdout data was documented before the test. Later ranking studies on separate development cohorts did not establish a reliable improvement. An audit of public real-world evidence found missing provenance and point-in-time fields; consequently no external performance evaluation was conducted. The findings characterize a synthetic evidence model, not operational SOC validity.
\end{abstract}
